% Pursuing Stacks (cotes 134-2 to 134-8): archive cote 134-2, site page 52 (PDF page 53).
% appendix (letters to Breen). Draft compared: ps1.tex:1270-1333.
% AI-transcribed by Claude Opus 5.5 (claude-opus-5-5, adaptive thinking, effort high), Message Batches API, batch ps1 (msgbatch_01Rv5hCaNTtUt4JHx1yqoHPL), 2026-11-12; not yet checked by eye.
\scanpage{134-2/52}
\mspage{\hand{37}} \mspage{5}

(en prolongeant \hand{$\mathcal{L}'$} par des 0 à droite). Le N- champ de Picard
strict correspondant n'est plus \hand{$\mathcal{F}$} mais \hand{$\mathcal{C}$}$^{N-n}$ \hand{$\mathcal{F}$} , où \hand{$\mathcal{C}$} est le
foncteur ``espace classifiant'', s'interprétant sur les n-catégories de Picard
strictes comme l'opération consistant à ``translater'' les i-objets en des
(i+1)-objets, et à rajouter un unique O-objet; il se prolonge aux n-champs \inserted{de Picard}
``de façon évidente'', on espère, de façon à commuter aux opérations d'image
inverse de n-champ\over{\illegible{1}}{s}. On aura donc pour i \hand{$\leq$} N
\[ \deleted{\illegible{1}}\hand{H}^{i}(X, \hand{\mathcal{L}'}) = \hand{\pi}_{N-i}\hand{\Gamma}_{X}(\hand{\mathcal{C}}^{N-n}\hand{\mathcal{F}}) \qquad (i \hand{\leq} N) \ . \]
Ceci posé, il s'impose, pour tout n-champ de Picard strict \hand{$\mathcal{F}$} sur X, de
poser
\[ \boxed{H^{i}(X , \hand{\mathcal{F}} ) = \hand{\pi}_{N-i}\hand{\Gamma}_{X}\hand{\mathcal{C}}^{N-n} \hand{\mathcal{F}} ) \qquad \text{si} \quad N \hand{\geq} i,n} \]
ce qui ne dépend pas du choix de l'entier N \hand{$\geq$} Sup(i,n) [\uline{NB} On a un
morphisme canonique de (n-1)- group\hand{\unsure{\"o}}ides ,
\[ \hand{\mathcal{C}}(\hand{\Gamma}_{X}\hand{\mathcal{F}}) \hand{\longrightarrow} \hand{\Gamma}_{X}(\hand{\mathcal{C}\mathcal{F}}) \ , \]
comme le montrent les constructions évidentes en termes de complexes de
cochaines, et on voit de même que celui-ci induit des isomorphismes pour
les \hand{$\pi$}$_{i}$ pour 1 \hand{$\leq$}i \hand{$\leq$} n+1 .]\hand{\unsure{$_4$|,}}

\uline{NB} On voit en passant que pour un n-champ en groupoides \hand{$\mathcal{F}$} sur X, si
on se borne à vouloir définir les H$^{i}$ (X, \hand{$\mathcal{F}$}) pour 0 \hand{$\leq$} i\hand{$\leq$}n, on n'a pas
besoin sur \hand{$\mathcal{F}$} d'une structure de Picard, car \over{\illegible{1}}{il} suffit de poser
\[ H^{i}( X , \hand{\mathcal{F}} ) = \hand{\pi}_{n-i}(\hand{\Gamma}_{X}(\hand{\mathcal{F}})) \qquad 0 \hand{\leq} i \hand{\leq} n. \]
Si d'autre part \hand{$\mathcal{F}$} est un n-Gr- champ (i.e. muni d'une loi de composition
\hand{$\mathcal{F}$} x \hand{$\mathcal{F}$} \hand{$\longrightarrow$} \hand{$\mathcal{F}$} ayant les propriétés formelles d'une loi de groupe)
\deleted{(}le (n+1) ``champ classifiant''\deleted{)} est défini, et on peut définir H$^{i}$(X , \hand{$\mathcal{F}$} )
pour \over{\illegible{1}}{i} $\leq$ n+1 par
\[ H^{i}( X , \hand{\mathcal{F}} ) = \hand{\pi}_{n+1-i} ( \hand{\Gamma}_{X}(\hand{\mathcal{C}\mathcal{F}})) \]
en particulier
\[ H^{n+1} (X , \hand{\mathcal{F}}) = \hand{\pi}_{o}(\hand{\Gamma}_{X}(\hand{\mathcal{C}\mathcal{F}})) = \text{sections de } \hand{\mathcal{C}\mathcal{F}} \text{ à équiva-}\]
lence\deleted{s} près .

\mspage{\hand{52}}
